\paragraph{Colour strings and fixed-source spatial records.}
For an open chain with zero external colour flux, Gauss elimination gives
$E_x=\sum_{y\leq x}C_y$ and a positive electric action
$\kappa\sum_x E_x^2$. The native three-quark singlet identity
$C_X^2=2n_X(3-n_X)/3$ implies the exact string cost
$4\kappa(x_{\max}-x_{\min})/3$.
On two cells, replacing $U$ by $U-\kappa/2$ preserves the entire
colour-singlet strong and colour-preserving contact action. Integer
checks retain all $120$ baryon and $2{,}520$ nuclear singlets.
With the same coefficients on chains of two through five cells, the
complete singlet dimensions are $120$, $364$, $816$ and $1{,}540$.
The neutral lowest energies are approximately
$0.93956542194$, $0.9396642972$, $0.939664299718$ and
$0.939664299718\,\mathrm{GeV}$. The two-cell mass is calibrated;
the larger-chain values are fixed-source outputs. Their axial response
is approximately $1.6658$, so the fitted two-cell axial value does not
transfer under this spatial extension. A positive source bound controls
relative-separation tails uniformly in chain length.
For one shared physical length, the exact relation
$\sigma^2Q=(4\kappa m_e/3)^2Q_{\rm graph}$ tests the string tension
and nuclear quadrupole together. Using the stated quadrupole comparison
and the illustrative $1\,\mathrm{GeV/fm}$ tension benchmark, every
$\kappa\leq2U$ preserving a nonnegative local colour cost fails the
simultaneous match. The resulting upper tension is below
$0.429\,\mathrm{GeV/fm}$. This is a finite-source comparison, without
a continuum confinement or string-breaking claim.
